%!TEX root = ../hutbthesis_main.tex

\chapter{绪论}

\section{研究背景及意义}

智能网联汽车，自动驾驶技术以及功能型无人装备正在飞速发展，以无人车为核心的部分智能化交通系统已在园区接驳，物流配送，安防巡检，港口作业，城市公交等多种场景达成大规模落地应用。无人车的环境感知，路径规划，行为决策，运动控制以及人车协同交互能力，会直接左右其行驶的安全性，通行的效率以及社会的接受程度，在无人车算法研发，功能验证和性能评价阶段，以往依靠实际道路开展考察的方式存在诸多明显问题，比如考察成本较高，危险场景不易重现，极端工况涵盖不够全面，研发历时过长等情况，这些状况已无法适应现代自动驾驶系统高效更新换代的需求，在此情形之下，数字化仿真考察由于具备成本低，安全性高，场景可按需定制，全过程可再现等特点，渐渐转变成无人车从算法设计到工程实施的关键支持手段\myrefcite{1}。

当下，主流的人车仿真平台包含CARLA，HUTB，Prescan以及Unreal
Engine等，这些平台有着比较完备的车辆动力学模型，交通流模拟系统，多种传感器仿真功能，环境渲染能力和物理碰撞处理能力，可以支持基本的无人车行驶检测，感知算法校验以及交通场景模拟。但在实际工程应用的时候，其一，平台控制很大程度上依靠代码级别的API调用，并且需要人工设置参数，操作流程很繁杂，对于非开发者来说入门难度较大。其二，场景创建，车辆控制，天气设定，传感器调用等功能模块彼此分离，缺少统一的交互通道和调度体系。其三，传统的仿真系统很难同大型模型，语音交互，Web可视化等智能方法很好地结合在一起，依靠自然语言激发的智能交互能力十分短缺。其四，环境设置过程极为繁琐，而且经常会出现依赖关系矛盾的情况，服务启动时的步骤也很复杂，轻量级部署以及快速重复利用的能力比较弱，这就不利于科研和工程项目尽快得到验证。

模型上下文协议（Model Context
Protocol，MCP）被提出之后，就给AI大模型同外部工具，硬件系统以及仿真平台之间的标准化通信带来了轻量化的解决办法，FastMCP是一个高性能，轻量级而且容易做二次开发的MCP实现框架，它具有低时延，高稳定性，便于封装工具，跨平台适配等特性，可以将人车仿真平台里零散分布的控制接口统一整合成标准的服务工具，从而做到从自然语言指令直接转为仿真执行指令，凭借前面这些技术基础，本文就设计并制作出依靠FastMCP的人车仿真器交互控制系统，用标准化协议来破除AI和仿真平台之间的障碍，创建起包含车辆控制，无人机控制，场景编辑，天气调节，传感器管理，Unreal引擎编辑等功能的智能交互系统，这对于优化人车仿真平台的易用性，智能化程度以及工程化效率有着重要的理论意义和实际价值。

在面向人车协同仿真的交互控制系统中，国内相关工作大多集中在驾驶模拟器的硬件设计、软件实现以及具体的交互场景上。早期主要是搭建真实感强的驾驶模拟平台，利用六自由度运动平台、高分辨率视景系统及力反馈方向盘等硬件，给用户带来身临其境的驾驶体验。这些系统对汽车动力学仿真以及驾驶员的行为研究都起到积极作用，是研究驾驶员在不同道路上行驶时如何操作必不可少的基础工具。但是传统的模拟器大多是基于特定需求进行开发，内部的核心仿真部分（例如汽车动力学模型、交通流模型）和交互控制紧密相连，因此很难对其进行扩展，也不方便加入新的交互设备或者满足不同的实验需求。

本文所介绍的``基于Fast
MCP的人车仿真器交互控制系统''，主要是在``快速''的基础上结合了``模型上下文协议''，为满足人车仿真实现了一套高效的交互控制协议框架。``Fast''并不只是指通信上的低延时，更重要的是这个协议栈可以支持仿真需要的确定性、及时性和大量数据传输能力。该系统的目标就是利用这个协议来实现高质量的车辆动力学模型、精确的环境感知模型以及丰富的人机交互界面（比如方向盘、油门、刹车甚至脑机接口）之间的良好连接，形成一个完整的交互控制系统。在这个系统中，所有的控制命令、传感器值、环境状况等都能以一种标准而且高效的方式在整个仿真过程中流动，进而使得驾驶员的动作、汽车的反应以及外界的影响都可以被精确地、毫秒级地模拟出来。这就可以解决以前由于各个模块之间的接口不同、数据难以有效地进行序列化或者中间件过重造成的种种问题。

本文的研究意义主要体现在以下四方面：

第一，创建起协议化，标准化而且低耦合的仿真交互框架，依靠FastMCP来统一封装各类别的仿真控制接口，从而减小模块整合的复杂程度并优化系统的可扩展能力。

第二,达成自然语言激发的人车仿真全流程控制，支持语音，Web界面，AI助手等多种渠道的交互方式，突出减小仿真操作的门槛。

第三,给出轻量化，虚拟化的一键式部署方案，达成环境自动设置，依赖自动安装，服务自动开启的目的，从而优化研发和检测的效率。

第四,创建起针对无人车功能验证，人车交互行为分析，自动驾驶决策算法检测以及交通仿真场景库形成这些研究方向的仿真支撑平台，该平台需具备高效性，灵活性与可复用性。

\section{国内外研究现状}

无人车仿真测试技术是自动驾驶系统开发的关键支撑手段。吴泰邦\myrefcite{1}针对无人车测试需求，设计并实现了交通仿真后端服务系统，为自动驾驶仿真提供了高效的数据服务支持。丁俊锋\myrefcite{2}面向自主驾驶场景，研究了无人车高效导航方法，重点改进了非结构化道路与动态障碍物环境中的路径规划性能，为仿真系统中的车辆运动控制提供了可靠的算法基础。黄嘉旻等\myrefcite{11}针对功能型无人车与通用智能网联汽车的差异，提出了一种面向功能验证的分层仿真场景库生成方法，构建了涵盖交通参与者交互、道路行驶、信号灯通行、交通标识识别、停车交互、障碍物识别的完整场景体系，并通过试验验证了该场景库的实用性和有效性。



在仿真系统性能优化方面，Fang等\myrefcite{25}利用分子动力学模拟辅助离子液体筛选，为复杂系统的仿真建模提供了方法学参考。Yan等\myrefcite{26}实现了天然丰度低$gamma$核素NMR信号的直接检测，其信号处理与数据分析方法对仿真系统中传感器数据的精确采集与处理具有借鉴意义。



人车交互行为建模是提升自动驾驶系统安全性和社会接受度的核心环节。胡馨宁\myrefcite{3}以无信号路段行人过街为研究重点，深入剖析了人车交互的行为特征和决策机理，构建了贴近实际交通行为的预测模型，有效提升了动态人车交互仿真的真实性。谢璞光等\myrefcite{4}提出了一种新型的人车交互实现方法，优化了车内外信息传达路径，加快了指令响应速度并增强了系统的易用性。曹哲\myrefcite{5}针对交叉口右转高冲突场景，构建了基于人车交互的自动驾驶车辆行为决策模型，提升了车辆在复杂冲突场景下的决策合理性和通行效率。赵彬越\myrefcite{6}在自动驾驶与人工驾驶共存的交通环境中，针对无信控路段的人车交互行为开展了系统仿真研究，揭示了交通流状态、运动参数对交互安全性和通行效率的影响规律。在国际研究方面，Yang等\myrefcite{16}利用多层路径分析法，全面解析了共享空间内人车交互的主要影响因素，为行车交互安全设计提供了量化依据。Li等\myrefcite{17}采用基于智能体（Agent）的方法构建了无信号路段人行横道上的人车交互模型，显著改善了动态交通场景仿真的精确程度和可信水平。Muduli等\myrefcite{18}依托长短期记忆网络（LSTM）构建了混合交通环境下的人车交互风险等级预测模型，为仿真测试中危险场景的自动生成及安全评价提供了算法支持。Lu等\myrefcite{23}针对智能座舱提出了交互增强方案，通过识别驾驶意图改善了人车交互体验，提升了座舱系统的智能化水平和流畅性能。张志杰等\myrefcite{12}从人机交互设计视角出发，按照人车交互原理提出了自动驾驶汽车多模态外显交互界面的设计方案，利用可视化提示增强了行人对车辆行驶意图的认知，进而改善了交互的安全性和连贯性。在交通系统优化方面，Maleki等\myrefcite{19}针对多分配枢纽最大覆盖问题提出了高效模型，其优化方法可应用于仿真系统中多智能体的路径规划与资源调度。徐婷婷等\myrefcite{9}研究了"车-路-网"耦合下含多类型充电桩协同的充电站规划，为智能网联交通系统的仿真建模提供了基础设施层面的理论支撑。

数字孪生和虚实融合技术为仿真系统带来了更高的真实度和交互能力，成为近年来人车仿真领域的重要发展方向。刘君丽\myrefcite{7}针对自动驾驶场景构建了数字孪生信息交互系统，实现了物理实体与虚拟仿真空间的数据同步、状态映射以及实时指令交互，改善了仿真系统与现实世界的一致性。李嘉然等\myrefcite{8}依照数字孪生理念设计了无轨胶轮车运动仿真与人机交互系统，并在工业特种车辆场景中验证了交互控制的即时性、稳定性和工程应用价值。

在底层技术支撑方面，高精度时间同步和低延迟通信是保障数字孪生系统实时性的关键。Ma等\myrefcite{24,30}针对快速定时探测器的性能优化开展了深入研究，探讨了微通道板光电倍增管信号链路的设计方法，为皮秒级时间分辨率系统的构建提供了技术参考。Lee等\myrefcite{27}提出了一种基于FPGA的时间数字转换器设计，采用反向梯度传播和虚拟箱方法减少时间箱之间的相关性，在资源受限的情况下实现了高分辨率和高线性度，已应用于实际快速定时测量中。这些研究成果对于设计低延迟、高确定性的交互控制协议底层通信机制具有重要指导意义。

在复杂系统建模优化方面，代理模型技术的应用为提升仿真效率提供了新思路。李良星等\myrefcite{10}采用物理信息神经网络建立了用于预测流场的代理模型，在保证一定准确性的前提下将预测速度提高了数个数量级，并大幅减少了内存占用。这一方法对于人车仿真中计算量巨大的模块（如复杂外部环境流动或精细传感器模拟）具有借鉴意义，可通过FastMCP协议将其接入交互控制环路，在保证仿真真实性的同时满足实时性要求。

在材料科学与仿真交叉领域，Nyemann等\myrefcite{20}研究了LiF:Mg,Cu,P的复合寿命特性，其脉冲光激发发光测量技术为仿真系统中光学传感器建模提供了物理参数参考。Pathak等\myrefcite{21}探讨了DPP-4抑制剂对脂肪肝病变的影响，其长期效应评估方法对仿真系统中车辆长期运行状态的稳定性分析具有方法论启示。Bauer-Fasching等\myrefcite{15}在Zwicky瞬变设施巡天中发现了新的ACV变量，其大规模天文观测数据处理方法对仿真系统中海量传感器数据的实时处理与分析具有参考价值。Bowles等\myrefcite{14}提出了从MS2时间序列数据中推断转录动力学参数的可扩展方法，其时间序列分析技术可应用于仿真系统中车辆轨迹数据的建模与预测。随着大模型技术的发展，自然语言交互成为降低仿真系统使用门槛的重要途径。Zhang等\myrefcite{13}对多无人机辅助反向散射通信系统进行了性能分析，为多智能体系统的协同控制提供了理论基础。Van
Dam\myrefcite{28}提出了基于多模态GUI架构的LLM对话助手接口方案，为自然语言与仿真系统的交互设计提供了架构参考。

在地球系统模拟与气候预测方面，Planton等\myrefcite{29}研究了海洋热含量对ENSO可预测性的不对称影响，其耦合环流模型中的预测方法对仿真系统中复杂环境参数的动态调控具有参考价值。然而，当前将自然语言交互与专业仿真平台深度结合的研究仍较为缺乏，特别是在自动驾驶仿真领域，如何实现从自然语言指令到仿真执行的高效、准确映射仍是亟待解决的技术难题。



综合国内外相关研究来看，无人车仿真检测，人车交互行为建模，数字孪生交互系统，自动驾驶决策算法这些方面已有诸多研究成果，给本文系统的设奠定下稳固的理论及技术根基，不过当下的平台和方法还是存在一些不够之处。

仿真平台接口分散、模块独立，缺乏统一的标准化交互框架，难以与AI大模型高效协同；

操作方式以代码调用与手动配置为主，自然语言交互能力薄弱，使用门槛较高；

场景构建、车辆控制、环境配置流程长难，难以快速响应复杂测试与多场景验证需求；

系统部署很复杂，环境设置也很麻烦，缺乏轻量化，一站式而且可以快速复用的部署方案。

依托FastMCP，本文设计出人车仿真器交互控制系统，该系统采用分层解耦架构，并将标准化工具加以封装，具备多渠道智能交互能力，而且经过轻量化一键部署就能运行，这样就填补了当前仿真平台存在的漏洞，全方位加强人车仿真系统的智能化，便捷化以及工程化水平。

\section{研究内容与技术路线}

\subsection{主要研究内容}

系统需求及架构设计方面，要剖析人车仿真控制，自然语言交互，轻量化部署这些关键需求，并规划依靠FastMCP的三层架构，该架构包含用户交互层，MCP服务层以及仿真器层，需明确各层的功能，模块划分情况和接口规范。

依靠HUTB/CARLA仿真平台来封装车辆，天气，传感器以及场景编辑等方面的工具接口。凭借Unreal
Engine开发编辑器，蓝图，节点以及UI编辑等相关的控制工具，从而创建起完整的工具链。

开发Web可视化交互界面，该界面可执行指令输入，状态反馈以及日志输出功能。利用Python虚拟环境做到依赖隔离，并编写一键启动脚本，从而达成环境设置，服务启动的全面自动化。

系统整合及检测验证包括MCP服务、交互界面以及仿真平台的联合调试，在一般仿真环境中进行功能测试与压力测试以检查系统对指令的反应是否正常、控制是否稳定、能否适应不同的场景并且部署是否高效等。

\subsection{技术路线以及主要工作}

本文按照 ``需求分析→架构设计→模块开发→系统集成→测试优化''
的完整技术路线开展研究：

需求分析阶段：深入分析CARLA仿真平台的API调用痛点、环境调控需求及部署流程瓶颈，明确系统功能需求与性能指标，确立基于FastMCP的服务化封装方案。

架构设计阶段：设计``用户交互层---MCP服务层---仿真执行层''的三级分层架构。明确HUTB
MCP Server（负责CARLA仿真控制）与Unreal MCP
Server（负责引擎编辑扩展）的功能边界与通信协议，制定工具封装规范与接口标准。

模块开发阶段：基于FastMCP框架完成仿真工具集的开发，实现车辆控制、天气调节、传感器管理等功能模块；开发自然语言指令解析与参数映射模块；编写自动化部署脚本与虚拟环境配置方案。

系统集成阶段：完成MCP服务、Web交互界面、CARLA仿真平台的联合调试，实现从自然语言输入到仿真执行的完整闭环控制。

测试验证阶段：在标准城市场景中开展功能测试、性能测试与压力测试，验证系统的指令响应正确性、部署效率及长时间运行稳定性，总结系统优势与改进方向。

\subsection{本文框架结构}

本文就依托FastMCP的人车仿真器交互控制系统的规划与达成予以阐述，总体上依照``理论分析，需求分析，架构设计，模块完成，系统评定，总结展望''这一整套研究脉络。

第一章属于绪论部分，该部分要阐述自动驾驶以及人车交互仿真相关研究的背景及其意义，还要整理国内外有关无人车仿真，人车交互，数字孪生以及智能交互系统的研究情况，从中找出当下仿真平台存在的一些问题，比如操作比较复杂，接口比较分散，自然语言交互能力有所短缺，部署起来也很繁杂之类的情况，从而明确本文的研究范围，技术路径以及总体框架。

第二章是研究的理论基础部分，该部分着重阐述MCP协议，FastMCP通信框架，HUTB/CARLA仿真器的工作原理，UnrealEngine编辑器的扩展机制，还有自然语言指令向仿真API实施映射的逻辑，这些内容会给系统架构设计以及模块开发给予理论支持和技术参考。

第三章介绍系统主要实现的技术，它从整体架构、双服务并行机制、仿真控制、引擎扩展、自然语言处理、模块化工具以及轻量化部署等方面进行阐述，详细介绍了系统的相关技术和实现过程，包括从底层通信到上层交互的整体设计思路。

第四章是系统详细设计与达成的部分，要明确开发环境，针对整体框架，功能模块，核心服务，指令映射，Web界面以及一键部署流程展开细致设计，还要给出关键代码和运行流程，从而全面展示系统工程化的达成过程。

第五章是总结与展望部分，该部分要总结文章的研究成果以及系统的改进之处，也要表明当下的短缺，还要从场景丰富程度，多模态交互，算法改良，分布式安排，数字孪生结合等方面来展望后续的研究方向。
